\chapter{Tinea Versicolor (Pityriasis Versicolor)}
\index{Tinea Versicolor (Pityriasis Versicolor)}
\label{sec:Tinea Versicolor (Pityriasis Versicolor)}

\section{Overview}
Tinea versicolor, also known as pityriasis versicolor, is a common superficial fungal infection of the skin caused by the yeast Malassezia globosa (formerly Pityrosporum orbiculare). It is characterized by hypo- or hyperpigmented scaly patches on the trunk and proximal extremities. The condition is most prevalent in adolescents and young adults in warm, humid climates.
\section{Classification}

\begin{description}[font=\normalfont\bfseries,leftmargin=3em,labelwidth=2.5em]
  \item[Cause] Infectious (Fungal)
  \item[Organ System] Skin
  \item[Pathophysiology] Overgrowth of Malassezia yeast in its pathogenic mycelial form
  \item[Duration] Chronic (months to years, recurrent)
  \item[Onset] Gradual
  \item[Mode of Transmission] Endogenous (normal skin flora); not contagious
  \item[Clinical Course] Chronic, recurrent; highly responsive to treatment
  \item[Severity] Mild
  \item[Extent of Disease] Localized to Regional
  \item[Age Group] Adolescents and Young Adults
  \item[Epidemiology] Prevalence 1--50\% depending on climate; highest in tropical regions
  \item[ICD-11 Code] 1F2C.0
\end{description}

\section{Causes}
Tinea versicolor is caused by overgrowth of Malassezia, a lipophilic yeast that is part of normal skin flora. The yeast converts from its saprophytic yeast form to the pathogenic mycelial form under conditions of warmth, humidity, sweat, and sebum. Malassezia produces azelaic acid, which inhibits melanogenesis and causes the characteristic hypopigmentation. Hyperpigmented forms result from inflammation-induced melanocyte stimulation.

\section{Risk Factors}

\begin{itemize}[noitemsep]
  \item Warm, humid climate and excessive sweating
  \item Oily skin and seborrheic areas
  \item Adolescence and young adulthood
  \item Immunosuppression
  \item Corticosteroid use
  \item Malnutrition
  \item Genetic predisposition
  \item Occlusive clothing
\end{itemize}

\section{Signs and Symptoms}

\subsection{Early symptoms}
\begin{itemize}[noitemsep]
  \item Early manifestations of tinea versicolor (pityriasis versicolor) may be subtle and nonspecific
\end{itemize}

\subsection{Common symptoms}
\begin{itemize}[noitemsep]
  \item \subsection{Early symptoms}
  \item \subsection{Common symptoms}
  \item \textbf{Common symptoms:}
  \item Well-demarcated, round to oval macules and patches on trunk, back, chest, shoulders
  \item Hypopigmented (most common), pink, or hyperpigmented lesions
  \item Fine, branny scale that becomes apparent with gentle scraping (Besnier sign)
  \item Lesions are typically asymptomatic but may cause mild pruritus
  \item Lack of tanning in affected areas (cosmetic concern)
  \item Face and neck involvement in children
  \item Difficulty performing daily activities (cosmetic concern)
\end{itemize}

\subsection{Less common symptoms}
\begin{itemize}[noitemsep]
  \item Atypical or less common presentations of tinea versicolor (pityriasis versicolor) may occur
\end{itemize}

\subsection{Emergency warning signs}
\begin{itemize}[noitemsep]
  \item Severe or worsening symptoms of tinea versicolor (pityriasis versicolor) require immediate medical evaluation
\end{itemize}
\section{Progression and Stages}
Lesions develop gradually over weeks to months and may persist indefinitely without treatment. The condition tends to recur after treatment, especially in warm, humid conditions. Spontaneous resolution can occur in cooler, drier climates. The pigmentary changes may persist for weeks to months after the fungus has been eradicated.

\section{Complications}

\begin{itemize}[noitemsep]
  \item Persistent pigmentary changes and cosmetic concern
  \item Recurrence (common, up to 80\% within 2 years)
  \item Mild pruritus, particularly with sweating
  \item Extensive skin involvement in immunosuppressed patients
  \item Reduced quality of life
  \item Emotional and psychological effects
\end{itemize}

\section{Diagnosis}

\begin{itemize}[noitemsep]
  \item Clinical diagnosis based on characteristic appearance and distribution
  \item Potassium hydroxide (KOH) preparation showing \"spaghetti and meatballs\" appearance
  \item Wood lamp examination showing yellow-green fluorescence
  \item Fungal culture (rarely needed)
  \item Skin biopsy in atypical cases
\end{itemize}

\section{Prevention}

\begin{itemize}[noitemsep]
  \item Keep skin clean and dry, especially after sweating
  \item Wear breathable, loose-fitting clothing
  \item Avoid excessive heat and humidity when possible
  \item Prophylactic use of antifungal shampoo periodically in susceptible individuals
  \item Go for regular check-ups so any problems are caught early
\end{itemize}

\section{Treatment Options}

\begin{itemize}[noitemsep]
  \item Topical antifungal agents: ketoconazole 2\% shampoo (single application or weekly)
  \item Selenium sulfide 2.5\% lotion applied for 10 minutes daily for 7 days
  \item Topical azole creams (clotrimazole, miconazole, econazole) for limited disease
  \item Oral antifungal therapy (fluconazole 300 mg single dose, itraconazole 200 mg daily for 5--7 days) for extensive cases
  \item Terbinafine 1\% solution or cream
  \item Lifestyle changes and self-care
\end{itemize}

\section{Living with It}

\begin{itemize}[noitemsep]
  \item Follow your treatment plan as prescribed
  \item Keep up with regular appointments with your healthcare provider
  \item Adjust your daily habits as needed
  \item Reach out to family, friends, or support groups for help
  \item Keep learning about your condition and how to manage it
\end{itemize}

\section{Outlook}
Tinea versicolor responds well to treatment, with clinical improvement seen within 1--3 weeks. However, recurrence rates are high (60--80\% within 1--2 years) in susceptible individuals. Maintenance therapy with intermittent ketoconazole shampoo or selenium sulfide may reduce recurrence. The pigmentary changes may take months to resolve even after the infection is treated. The condition is benign and does not cause long-term health problems.
